\chapter{Introdução}
\label{chap:introducao}

Os mercados de derivativos permitem negociar a volatilidade como um ativo:
opções, \textit{swaps} de variância e índices de volatilidade implícita, como o
VIX, transformam a incerteza sobre o preço futuro em algo que se compra e se
vende \citep{carr2009variance, alexanderImeraj2019}. Nos criptoativos, essa
incerteza é muito maior que nos mercados tradicionais. Entre 2017 e 2022, a
volatilidade implícita média das opções de Bitcoin foi de 82\% a.a., muito
acima da do S\&P~500 \citep{almeida2024cryptoVRP}; na amostra desta
dissertação, de março de 2021 a março de 2026, foi de $61{,}9\%$ a.a. As
opções de Bitcoin se concentram na Deribit, a maior bolsa de derivativos de
criptoativos e a principal fonte de dados dos estudos sobre essas opções
\citep{almeida2024cryptoVRP}; em 2019 e 2020, quase toda a negociação de opções
de Bitcoin ocorria nela \citep{alexanderImeraj2019}. Desde 24 de março de 2021,
a Deribit publica o DVOL, índice de volatilidade implícita de 30 dias
construído a partir dos preços dessas opções \citep{han2026price}
(Seção~\ref{sec:refer-dvol}).

O preço da volatilidade nas opções não é uma previsão neutra. Em ações, a
volatilidade implícita supera, em média, a que se realiza depois: os
compradores de opções pagam um prêmio para se proteger do risco de
volatilidade, e esse prêmio prevê os retornos futuros do índice
\citep{bollerslev2009expected}. Esta dissertação estuda o prêmio de risco de
volatilidade do Bitcoin (BVRP), a diferença entre a volatilidade histórica (VH)
dos 30 dias seguintes e a volatilidade implícita (IV) de 30 dias dada pelo
DVOL, $\text{BVRP}_t = \text{VH}_{t+1:t+30} - \text{IV}_t$. Com essa convenção,
um prêmio maior corresponde a um BVRP mais negativo. A dissertação responde a
quatro perguntas: o prêmio existe? Ele pode ser previsto com a informação
disponível em $t$? O prêmio previsto prevê linearmente os retornos do Bitcoin?
Um modelo não linear encontra a relação que o linear não encontra?

\section*{Roteiro e resultados}
\label{sec:intro-roteiro}

\textbf{Existência do prêmio.} A dissertação mede o prêmio com a VH dos 30 dias
seguintes, e não com a dos 30 anteriores, porque a IV de $t$ precifica a
volatilidade de $t+1$ a $t+30$: só a janela prospectiva cobre o mesmo período.
A versão com a VH passada,
$\text{BVRP}^{\text{proxy}}_t = \text{VH}_{t-29:t} - \text{IV}_t$, é conhecida
em $t$ e serve de variável explicativa, mas supõe que a volatilidade siga um
passeio aleatório. A única hipótese formal da dissertação (H1) é que o prêmio é
negativo em média. O teste regride o BVRP numa constante, com erro-padrão de
Newey--West de 31 defasagens, que cobre a sobreposição entre alvos vizinhos.
Nas 1.806 datas da amostra, a média é de $-8{,}76$ p.p. ($t = -5{,}15$), e a IV
supera a VH dos 30 dias seguintes em 72\% das datas: o prêmio existe. A
distribuição da VH de 30 dias é bimodal, o que indica dois regimes de
volatilidade; a dissertação estima o corte entre eles, de $37{,}3\%$ a.a., só
com os dois primeiros anos da amostra, para que nenhuma data seja classificada
com informação posterior a ela. No regime de baixa volatilidade, a IV supera a VH
dos 30 dias anteriores, mas não a dos 30 seguintes: a proxy registra um prêmio
maior nesse regime, e o BVRP, nenhum.

\textbf{Previsão do prêmio.} O BVRP de $t$ só é conhecido em $t+30$ e não pode
ser usado como variável explicativa, e a proxy aproxima mal a sua realização: a
correlação entre as duas na mesma data é de $-0{,}02$. Por isso, antes de
relacionar o prêmio aos retornos, a dissertação prevê o BVRP fora da amostra,
com 17 variáveis conhecidas em $t$ --- volatilidades históricas de 1 a 90 dias,
a IV e suas variações, retornos passados e a proxy --- e cinco modelos: MQO,
Ridge, LASSO, floresta aleatória e XGBoost. A dissertação reestima os modelos a
cada 30 dias em janela expansiva, e o treino de cada origem exclui os alvos que
ainda não se realizaram. Na janela expansiva, todos superam a média histórica: o Ridge reduz o erro
quadrático em 22\%, e o teste de \citet{clark2007approximately} rejeita a
igualdade com a média para quatro dos cinco modelos, mesmo com a correção para
múltiplas comparações. A previsibilidade, porém, é modesta: as previsões ficam
próximas da média e não antecipam os episódios de prêmio positivo. A relação
entre as variáveis e o prêmio muda entre os regimes dentro da amostra, mas essa
mudança não melhora a previsão fora dela.

\textbf{Previsibilidade linear dos retornos.} Com o BVRP previsto, a
dissertação estima a regressão de \citet{bollerslev2009expected}: o retorno do
Bitcoin nos $h$ dias seguintes, com $h$ de 1 a 60 dias, sobre o prêmio em $t$.
O BVRP previsto é um regressor gerado, e o erro-padrão usual ignora a incerteza
da sua estimação \citep{pagan1984}; por isso, a inferência principal usa um
bootstrap em blocos em dois níveis, que reestima o modelo de previsão em cada
reamostragem. O coeficiente é negativo nos seis horizontes, na direção de
\citet{bollerslev2009expected}: um prêmio maior antecede retornos maiores.
Nenhum horizonte, porém, é significante depois das correções para o regressor
gerado e para as comparações múltiplas, e o teste conjunto nos seis horizontes
não rejeita ($p = 0{,}37$); ele só rejeita quando ignora a estimação do BVRP
previsto. A inclinação é mais negativa no regime de baixa volatilidade, mas a
diferença entre os regimes depende da posição do corte.

\textbf{Previsibilidade não linear.} Como a diferença entre os regimes depende
de um corte fixado de antemão, a dissertação repete o teste com uma floresta
aleatória, que escolhe sozinha os cortes nas volatilidades e gera o BVRP
previsto. O coeficiente fica perto de zero em todos os horizontes, e o teste
conjunto não rejeita ($p = 0{,}96$). As árvores cortam a VH de 30 dias entre
45\% e 58\% a.a., e não em $37{,}3\%$, e a floresta aplicada direto ao retorno
também não supera a média histórica de forma significante. O
Apêndice~\ref{app:estrategias} traduz o prêmio previsto em regras de negociação: nenhuma supera o compra e manutenção
em índice de Sharpe, e a perda máxima menor vem de ficar fora do mercado na
maior parte dos dias.

Em resumo, o prêmio existe e é grande, e as variáveis disponíveis em $t$
contêm alguma informação sobre ele, mas o prêmio previsto não prevê os retornos
do Bitcoin nesta amostra, nem com o modelo linear nem com o não linear. As 987
previsões, com alvos de 30 dias sobrepostos, equivalem a cerca de 33
observações independentes; os testes têm pouco poder, e a ausência de
significância não é evidência de ausência de relação.

\section*{Contribuição}
\label{sec:intro-contribuicao}

Dois trabalhos estudam o prêmio de volatilidade do Bitcoin com dados da
Deribit. \citet{alexanderImeraj2019} são os primeiros: constroem um índice de
volatilidade do Bitcoin pelo método do VIX, com mais de 7 milhões de preços de
opções coletados a cada 15 minutos entre março de 2019 e março de 2020, e medem
o prêmio pelo resultado de um \textit{swap} de variância, com a variância
realizada ao longo da vida do contrato --- uma definição prospectiva, com o
mesmo sinal da adotada aqui. Eles descrevem o prêmio em vencimentos de 7 a 90
dias, quase sempre negativo a partir de julho de 2019, e sua correlação com o
prêmio de outros ativos, próxima de zero antes da pandemia e mais alta depois,
com ações e ouro. \citet{almeida2024cryptoVRP} usam opções da Deribit de julho
de 2017 a dezembro de 2022 e medem o prêmio pela variância neutra ao risco
menos a física, com a variância realizada dos 27 dias anteriores --- uma medida
retrospectiva, com o sinal oposto ao desta dissertação. O prêmio médio, de 0,14
em variância anualizada, é muito maior que o do S\&P~500. Com dois regimes
obtidos pelo agrupamento das densidades neutras ao risco na amostra inteira,
eles encontram prêmio maior no regime de baixa volatilidade. Nenhum dos dois
prevê o prêmio nem regride o retorno do Bitcoin nele. Em trabalho relacionado,
\citet{han2026price} usa o DVOL diário desde o início da série, agregado por
semana: as inovações do índice, resíduos de um AR(1), formam um fator de risco
de volatilidade do mercado de criptoativos, e a sensibilidade de cada
criptoativo a esse fator é precificada no corte transversal dos retornos da
semana seguinte, de dezembro de 2021 a julho de 2024. Ele não mede o prêmio
de volatilidade.

Esta dissertação acrescenta quatro elementos. Primeiro, a amostra: cinco anos
de dados diários, de março de 2021 a março de 2026, mais longa que a de
\citet{alexanderImeraj2019} e mais recente que a de
\citet{almeida2024cryptoVRP}, porque cobre 2023 a 2026 e a aprovação dos ETFs
de Bitcoin à vista, em janeiro de 2024. A volatilidade implícita é o mesmo
índice usado por \citet{han2026price}, o DVOL, e não um índice próprio.
Segundo, a definição: o prêmio é prospectivo, como em
\citet{alexanderImeraj2019}, mas em volatilidade, e a proxy retrospectiva entra
como variável explicativa, e não como o prêmio. Terceiro, a previsão: a
dissertação prevê o prêmio fora da amostra e usa a previsão na regressão de
\citet{bollerslev2009expected}, com a inferência corrigida para o regressor
gerado. Quarto, os regimes: como \citet{almeida2024cryptoVRP}, a dissertação
usa dois regimes de volatilidade, mas definidos por um corte na VH estimado sem
informação posterior às datas classificadas. Com a proxy, retrospectiva como a
medida deles, o prêmio é maior no regime de baixa volatilidade --- média de
$-13{,}56$ p.p., contra $-7{,}26$ p.p. no de alta, lembrando que prêmio maior é
BVRP mais negativo ---, o mesmo padrão que eles encontram; com o BVRP
prospectivo, a diferença se inverte (Seção~\ref{sec:dados-regimes}).

\section*{Organização}
\label{sec:intro-organizacao}

O Capítulo~\ref{chap:referencial} apresenta o referencial teórico: a
volatilidade implícita e a histórica, a definição do prêmio, o mercado de
opções de Bitcoin e os modelos de aprendizado de máquina. O
Capítulo~\ref{chap:dados} descreve os dados, testa H1 e define os regimes. O
Capítulo~\ref{chap:metodologia} formaliza as questões, as amostras, os modelos
e a inferência. O Capítulo~\ref{chap:previsao_bvrp} prevê o BVRP, e os
Capítulos~\ref{chap:retornos_linear} e~\ref{chap:retornos_nao_linear} testam se
o BVRP previsto prevê os retornos, com o modelo linear e com o modelo em
árvore. O Capítulo~\ref{chap:consideracoes_finais} conclui, e a Nota
Metodológica trata da reprodutibilidade. O Apêndice~\ref{app:estrategias}
avalia as regras de negociação, e os Apêndices~\ref{app:variaveis},
\ref{app:scripts} e~\ref{app:pastas} documentam as variáveis, os scripts e a
estrutura de pastas.
